\def\Ponderation{37,5}
\def\SigleNumero{STT-1000}
\def\NomCours{Probabilités et statistique}
\def\DateEvaluation{18 décembre 2023}
\def\HeureEvaluation{18h30 à 21h30}
\def\Session{Automne 2022}
\def\NomEvaluation{Examen 2}

\def\OptionsExam{11pt}

% Ne rien mettre entre les accolades si l'enseignant (ou l'enseignante) est aussi le coordonnateur (ou la coordonnatrice)
\def\Coordonnatrice{}
\def\Coordonnateur{}

% Ne rien mettre entre les accolades s'il y a plus d'une section
\def\Enseignant{}
\def\Enseignante{}

% Ne rien mettre entre les accolades s'il s'agit d'un cours à section unique
% Mettre le nom de chacun des enseignants et enseignantes entre les accolades
% Une case à cocher sera générée pour chacune des sections
\def\SectionA{Ahmed Sid-Ali (en ligne)}
\def\SectionB{Julien Miron (en présence)}
\def\SectionC{}
\def\SectionD{}

% Concerne seulement les travaux pratiques : écrire entre les accolades le nombre maximal de membres pour une équipe
\def\NombreLignes{}		

% Si vous souhaitez qu'apparaisse une déclaration d'honneur au bas de la première page, écrivez quelque chose entre les accolades
\def\EnLigne{}

\newcommand{\affichepoints}{}     % ← AVEC points
% \renewcommand{\affichepoints}{on} % ← SANS points (non vide)

\def\Directives{
\begin{enumerate}
\item Matériel autorisé:
\begin{itemize}
\item Calculatrice scientifique {\bf autorisée par la faculté des sciences et de génie};
\item Les feuillets de formules fournis avec l'examen. Ceux-ci incluent les tables de loi nécessaires. Veuillez \textbf{ne rien écrire} sur les feuillets de formules.
\end{itemize}
\item Assurez-vous que cette évaluation comporte \numquestions ~questions réparties sur \numpages{}~pages, pour un total de \numpoints{}~points.\
\item Lorsque des boîtes de réponses sont présentes, vous \textbf{\textit{devez}} y écrire vos réponses.\
\item Lorsque nécessaire, utilisez deux chiffres après la virgule.\
\item \textbf{Veuillez ne pas dégrafer votre examen.}
\item Ajoutez vos initiales dans le bas de chaque page de l'examen.
\item Veuillez déposer votre carte d'identité admise avec photo sur votre bureau.\
\item Vous devez justifier votre réponse par une démarche claire, sauf pour les questions à choix de réponse.\
\item Le verso des feuilles peut être utilisé comme brouillon, \textbf{mais ne sera pas corrigé}. Attention à ne rien dégrafer.
\end{enumerate}
}

% Veuillez indiquer le type d'évaluation entre crochets
% Examen : option à choisir pour un examen présentiel qui sera corrigé de manière classique. Un tableau des points sera affiché sur la page couverture.
% Gradescope : option à choisir pour un examen qui sera corrigé via Gradescope. Aucun tableau de pointage n'apparaîtra et il y a quelques différences au niveau de la mise en page
% TPS :
% TPF :
% Devoir :
\providecommand{\Gradescope}{}

\documentclass{Evaluation}
\usepackage{multicol,graphicx}

\printanswers
\begin{document}
\begin{questions}

\pageblanche
\question\label{QCMicTests}
Un intervalle de confiance de niveau 95\% pour l'espérance $\mu$ est $[ 5,11 ; 9,11]$. Pour chaque énoncé qui suit, indiquez si l'énoncé est vrai, faux ou si nous n'avons pas assez d'information pour le savoir.

\medskip

\begin{parts}
    \part[1] Au seuil 5\%, on ne rejette pas $H_0: \mu=5$ pour $H_1:\mu \neq 5$.

\medskip
    \begin{oneparcheckboxes}
        \CorrectChoice Vrai 
        \choice Faux
        \choice Nous n'avons pas assez d'information pour le savoir.
    \end{oneparcheckboxes}
\medskip
 
   \part[1] Au seuil 1\%, on rejette $H_0: \mu=4$ pour $H_1:\mu \neq 4$.

 \medskip  \begin{oneparcheckboxes}
        \choice Vrai 
        \choice Faux
        \CorrectChoice Nous n'avons pas assez d'information pour le savoir.
\medskip    \end{oneparcheckboxes}
    
   \part[1] Au seuil 10\%, on rejette $H_0: \mu=4$ pour $H_1:\mu \neq 4$.

\medskip   \begin{oneparcheckboxes}
        \CorrectChoice Vrai 
        \choice Faux
        \choice Nous n'avons pas assez d'information pour le savoir.
\medskip    \end{oneparcheckboxes}
    
   \part[1] Au seuil 5\%, on rejette $H_0: \sigma=1$ pour $H_1:\sigma \neq 1$.

\medskip   \begin{oneparcheckboxes}
        \choice Vrai 
        \CorrectChoice Faux
        \choice Nous n'avons pas assez d'information pour le savoir.
\medskip    \end{oneparcheckboxes}
    
%   \part[1] Au seuil 1\%, on rejette $H_0: \mu=8$ pour $H_1:\mu \neq 8$.

%\medskip   \begin{oneparcheckboxes}
%        \choice Vrai 
%        \CorrectChoice Faux
%        \choice Nous n'avons pas assez d'information pour le savoir.
%    \end{oneparcheckboxes}\medskip

\end{parts}

\vspace{0.5cm}
\question[4] Soit une loi normale d'espérance $\mu$ et de variance $\sigma^2=64$.
Nous voulons faire un test d'hypothèses au seuil $\alpha=0,05$ pour confronter 
\begin{align*}
    H_0:&\mu=600,\\
    H_1:&\mu\neq 600.
\end{align*} 
Les valeurs critiques du test sont donc $-z_{0,025}=-1,96$ et $z_{0,025}=1,96$.
Avec un échantillon aléatoire de $n=241$ étudiants, nous obtenons un seuil observé de $0,004$. 
Cochez tous les énoncés qui sont vrais parmi les suivants. \textit{\textbf{Une réponse cochée alors qu'elle ne devrait pas l'être entraîne une pénalité de 1 point. Le minimum possible pour cette question est 0 point.}}
 
\begin{checkboxes}
\correctchoice Nous savons que $|z_0|>1,96$.
\choice Nous savons que $|z_0|<1,96$.
\correctchoice Nous rejetons $H_0$.
\choice Nous ne rejetons pas $H_0$.
\choice Nous sommes certains que l'intervalle de confiance à 95\% pour $\mu$ contient 600.
\correctchoice Nous sommes certains que l'intervalle de confiance à 95\% pour $\mu$ ne contient pas 600.
\choice Nous sommes certains que l'intervalle de confiance à 99\% pour $\mu$
 contient 600.
\choice Nous sommes certains que l'intervalle de confiance à 99\% pour $\mu$ ne contient pas 600.
\choice Nous ne pouvons rien dire sur ce test car nous ne connaissons pas la statistique observée, $z_0$.
\end{checkboxes}
\vspace{0.5cm}


\question[4]
Parmi les statistiques descriptives suivantes, lesquelles sont des mesures de dispersion? \textit{\textbf{Une réponse cochée alors qu'elle ne devrait pas l'être entraîne une pénalité de 1 points. Le minimum possible pour cette question est 0 point.}}

\begin{checkboxes}
\correctchoice La moyenne échantillonnale
\choice La médiane échantillonnale
\choice Le mode échantillonnal
\correctchoice La variance échantillonnale
\choice L'écart interquartile
\choice L'étendue
\end{checkboxes}

\pageblanche
\question[4]
Soit $Z_1,Z_2$ et $Z_{3}$, des variables aléatoires normales centrées-réduites indépendantes et $\displaystyle W=Z_1^2+2Z_2^2+3Z_3^2$. 

Que vaut Var$(W)$? \textbf{\textit{Donnez les détails de votre démarche. Une réponse sans démarche ne donnera aucun point.}}

\begin{solution}
On a $Z_i^2 \sim \chi^2(1)$ pour $i=1,2,3$, donc $E[Z_i^2]=1$ et $\text{Var}(Z_i^2)=2$.

Comme les $Z_i$ sont indépendantes :
\begin{align*}
\text{Var}(W) &= \text{Var}(Z_1^2) + 4\,\text{Var}(Z_2^2) + 9\,\text{Var}(Z_3^2) \\
&= 2 + 4\times 2 + 9\times 2 \\
&= 2 + 8 + 18 = 28.
\end{align*}
\end{solution}

\vspace{9cm}
Var$(W)=$\boite{3}{1.5} 
\question 
Annabelle suppose que le temps d'attente dans un bureau de la SAAQ suit une loi exponentielle de moyenne 240 minutes.

Elle veut construire un intervalle de confiance bilatéral à $95\%$ pour le vérifier. 
Elle note alors le temps d'attente de $100$ personnes choisies au hasard et obtient une moyenne $\Bar{x}=250$, et une variance $s^2=59~536$. 
À partir de ces résultats, elle construit un intervalle de confiance en utilisant le quantile $z_{0,025}$, à partir de sa table de la loi normale, et obtient l'intervalle $[130,44;369,56]$. 
\begin{parts}
    \part[2] Écrivez les hypothèses du test qu'Annabelle devrait faire à partir de cet intervalle de confiance, en lien avec son interrogation.
    \begin{solution}
    La loi exponentielle de moyenne 240 a pour espérance $\mu = 240$. Le test associé à l'intervalle à 95\% est
    \[
        H_0: \mu = 240 \quad \text{contre} \quad H_1: \mu \neq 240.
    \]
    \end{solution}
    \vspace{2cm}
    \part[2] Quelle serait la décision du test en (a)?
    \begin{solution}
    L'IC à 95\% est $[130{,}44\,;\, 369{,}56]$. Puisque $240 \in [130{,}44\,;\, 369{,}56]$, on \textbf{ne rejette pas} $H_0$ au seuil 5\%.
    \end{solution}
    \vspace{2cm}
    \part[2] Le seuil du test sera-t-il exact? Pourquoi?
    \begin{solution}
    Non. Les temps d'attente suivent une loi exponentielle (non normale). L'IC a été construit avec le quantile $z_{0{,}025}$ grâce au TCL; le seuil est donc \textbf{approximatif}.
    \end{solution}
\end{parts}

\pageblanche
\question\label{rls}

Afin d'estimer le volume de bois disponible dans la forêt qui mène à La Barillette, nous aimerions modéliser le volume de bois d'un arbre à l'aide de sa circonférence. Nous disposons du volume ($Y$) et de la circonférence ($X$) de 31 arbres. Afin d'étudier le modèle $Y_i = \beta_0 + \beta_1 X_i + E_i$, en supposant que les postulats de régression linéaire sont respectés, nous avons résultats suivants:

\begin{align*}
	\sum_{i=1}^{31} x_i = 930;  \quad \sum_{i=1}^{31} x_i^2 = 28\;900 ; \quad \sum_{i=1}^{31} x_iy_i =113\;880;\quad
	\sum_{i=1}^{31} y_i = 3\;596; \quad \sum_{i=1}^{31} y_i^2 = 446\;716.
\end{align*}

\begin{parts}
	\part[8] Donnez les estimations de $\beta_0$ et $\beta_1$. \textbf{\textit{Donnez les détails de votre démarche. Des réponses sans démarche ne donneront aucun point.}}
	
	
	\begin{solution} $$ \widehat{\beta}_1 
		= \frac{S_{XY}}{S_{XX}}
		= \frac{\sum x_iy_i - n\overline{x}\overline{y}}{\sum x_i^2 - n\overline{x}^2} 
		= \frac{\displaystyle 13\; 887.86 -  \frac{410.7 \times  935.3}{31}}{\displaystyle 5\; 736.55 - \frac{410.7^2 }{ 31} }
		= \frac{1496.644}{295.4374}
		= 5.065858 $$
	
	$$ \widehat{\beta}_0 =  \overline{y}-\widehat{\beta}_1 \overline{x} 
		= \frac{935.3}{31} -  5.065858 \frac{410.7}{31} 
		= -36.94348 $$ 
		
		$$ \widehat{Y}_i = -36.94348 + 5.065858  x_i$$
		\end{solution}
	\vspace{17cm}

$\widehat{\beta}_0=$\boite{3}{1.5} $\widehat{\beta}_1=$\boite{3}{1.5}

\vspace{0.5cm}
	\pagesuivante{\ref{rls}}

\newpage
\suite{\ref{rls}}


	\part[2] Sophie aimerait estimer le volume de bois de l'arbre devant sa maison. Cet arbre a une circonférence de 28. La prévision ponctuelle du volume de bois de cet arbre est de 104. Choisissez l'intervalle de valeurs ayant 95\% de chances de contenir le volume de bois de cet arbre de circonférence 28.
	
	\begin{checkboxes}
\choice $ \displaystyle	104 \pm t_{0,05;29} \sqrt{MS_E \left[ \frac{1}{n} + \frac{(28-\overline{x})^2}{S_{XX}}\right]} $ 
\choice $ \displaystyle	104 \pm t_{0,025;1} \sqrt{MS_E \left[ \frac{1}{n} + \frac{(28-\overline{x})^2}{S_{XX}}\right]} $ 
\choice $ \displaystyle	104 \pm t_{0,025;29} \sqrt{MS_E \left[ \frac{1}{n} + \frac{(28-\overline{x})^2}{S_{XX}}\right]} $ \vspace{0.2cm}
\correctchoice $ \displaystyle		104 \pm t_{0,025;29} \sqrt{MS_E \left[ 1+ \frac{1}{n} + \frac{(28-\overline{x})^2}{S_{XX}}\right]} $
\choice $ \displaystyle		104 \pm t_{0,025;1} \sqrt{MS_E \left[ 1+ \frac{1}{n} + \frac{(28-\overline{x})^2}{S_{XX}}\right]} $
\choice $ \displaystyle		104 \pm t_{0,05;29} \sqrt{MS_E \left[ 1+ \frac{1}{n} + \frac{(28-\overline{x})^2}{S_{XX}}\right]} $
	\end{checkboxes}
\vspace{0.5cm}
		\part[1] Si Sophie avait un arbre ayant une circonférence de 35, l'intervalle de valeurs calculé à la question précédente serait... 
	\begin{checkboxes}
			\choice plus court 
			\choice de même longueur  
			\correctchoice plus long   
		\choice on ne peut pas savoir la différence de longueur 
	\end{checkboxes}
\vspace{0.5cm}

\part[1] La valeur-p du test de significativité du lien linéaire entre $X$ et $Y$ est égale à $0,00014$. Cochez l'énoncé qui est vrai.

	\begin{checkboxes}
            \choice Il n'y a pas de lien linéaire significatif entre $X$ et $Y$.
            \choice Il y a un lien linéaire significatif entre $X$ et $Y$.
            \correctchoice Nous ne pouvons rien dire car nous ne connaissons pas le seuil du test.
            \choice Nous ne pouvons rien dire car nous ne connaissons pas $t_0$.
	\end{checkboxes}

\end{parts}

\pageblanche
\question Soit $X_1,\hdots, X_{100}$, des variables aléatoires indépendantes et identiquement distribuées selon une loi normale d’espérance $\mu$ et de variance $\sigma^2 = 400$. 
Nous désirons construire un intervalle de confiance de niveau $95\%$ pour $\mu$. Bien sûr, l'intervalle de confiance aura la forme d'un des intervalles vus au cours. Nous notons cet intervalle $[C_1;C_2]$.
\begin{parts}
    \part[4] Que vaut $\mathbb{E}[C_2]$? \textbf{\textit{Donnez les détails de votre démarche.}}
    \begin{solution}
    On a $C_2 = \overline{X} + 1{,}96\,\dfrac{\sigma}{\sqrt{n}} = \overline{X} + 1{,}96\,\dfrac{20}{10} = \overline{X} + 3{,}92$.
    \[
        \mathbb{E}[C_2] = \mathbb{E}[\overline{X}] + 3{,}92 = \mu + 3{,}92.
    \]
    \end{solution}
    \vspace{4.5cm}
    \part[4] Que vaut $\text{Var}[C_2]$? \textbf{\textit{Donnez les détails de votre démarche.}}
    \begin{solution}
    \[
        \text{Var}(C_2) = \text{Var}(\overline{X}) = \frac{\sigma^2}{n} = \frac{400}{100} = 4.
    \]
    \end{solution}
    \vspace{4.5cm}
    \part[2] Quelle est la loi de probabilité de $C_2$?
    \begin{solution}
    Puisque les $X_i$ sont normales, $\overline{X}$ est normale, et $C_2 = \overline{X}+3{,}92$ l'est aussi:
    \[
        C_2 \sim \mathcal{N}(\mu + 3{,}92\;; 4).
    \]
    \end{solution}
    \vspace{2cm}
    \part[6]Quelle est la probabilité que $C_2$ soit inférieure ou égale à $\mu$? \textbf{\textit{Donnez les détails de votre démarche. Une réponse sans démarche ne donnera aucun point.}}
    \begin{solution}
    \begin{align*}
        P(C_2 \leq \mu) &= P\!\left(\overline{X} + 3{,}92 \leq \mu\right) \\
        &= P\!\left(\frac{\overline{X}-\mu}{2} \leq \frac{-3{,}92}{2}\right) \\
        &= P(Z \leq -1{,}96) \\
        &= 0{,}025.
    \end{align*}
    \end{solution}
    \vspace{8cm}

    Réponse: \boite{3}{1.5}
\end{parts}

\pageblanche
\question[10]
Soit les échantillons

\begin{align*}
Y_{1}, \ldots, Y_{n} \overset{i.i.d.}{\sim} \mathcal{N}(\mu_Y, \sigma^2_Y), & & X_{1}, \ldots, X_{m} \overset{i.i.d.}{\sim} \mathcal{N}(\mu_X, \sigma^2_X),
\end{align*}
où $n=8$, $m=8$ et $\sigma^2_Y=\sigma^2_X$ sont inconnues. Pour le test dont les hypothèses sont 
\begin{align*}
    H_0: \mu_Y-\mu_X=0,\\
    H_1: \mu_Y-\mu_X>0,
\end{align*}
 la valeur-p est de $0,01$. Sachant que $\bar{y}=6,824$, $\bar{x}=4,2$ et $s^2_Y=4,5$, quelle est la valeur de $s^2_X$? \textbf{\textit{Donnez les détails de votre démarche. Une réponse sans démarche ne donnera aucun point.}}
\begin{solution}
Test unilatéral droit avec $n+m-2=14$ degrés de liberté. La valeur-p est 0{,}01, donc $t_0 = t_{0{,}01\,;\,14} \approx 2{,}624$.

La statistique vaut
\[
t_0 = \frac{\bar{y}-\bar{x}}{s_p\sqrt{\frac{1}{n}+\frac{1}{m}}} = \frac{6{,}824-4{,}2}{s_p\,\sqrt{\frac{2}{8}}} = \frac{2{,}624}{s_p / 2}.
\]
Donc $s_p = \dfrac{2{,}624 \times 2}{2{,}624} = 2$, soit $s_p^2 = 4$.

Par la formule de la variance poolée :
\begin{align*}
s_p^2 &= \frac{(n-1)s_Y^2 + (m-1)s_X^2}{n+m-2} \\
4 &= \frac{7\times 4{,}5 + 7\,s_X^2}{14} \\
56 &= 31{,}5 + 7\,s_X^2 \\
s_X^2 &= \frac{24{,}5}{7} = 3{,}5.
\end{align*}
\end{solution}
\vspace{18cm}
$s^2_X=$ \boite{3}{1.5}

\pageblanche
\question[10] Soit $X_1,\hdots, X_{25}$, des variables aléatoires indépendantes et identiquement distribuées selon une loi normale d'espérance~$\mu$ et de variance $\sigma^2=625$.
Nous désirons effectuer un \textbf{test unilatéral à droite} avec $H_0:\mu=9$ au seuil $\alpha = 0,025$.
Nous avons les résultats suivants:
\begin{align*}
    \bar{x}= 11,34 &  & s^2=754.
\end{align*}
Si nous supposons qu'en réalité la vraie moyenne $\mu=\mu_1$ et que la puissance du test est $0,13567$, que vaut $\mu_1$? 
\textbf{\textit{Donnez les détails de votre démarche. Une réponse sans démarche ne donnera aucun point.}}
\begin{solution}
Avec $\sigma=25$ et $n=25$, on a $\sigma/\sqrt{n}=5$.
La région de rejet est $\overline{X} > 9 + z_{0{,}025}\times 5 = 9 + 1{,}96\times 5 = 18{,}8$.

La puissance est
\[
\text{Puiss}(\mu_1) = P\!\left(\overline{X}>18{,}8\mid \mu=\mu_1\right)
= P\!\left(Z > \frac{18{,}8-\mu_1}{5}\right) = 0{,}13567.
\]

Donc $P\!\left(Z \leq \dfrac{18{,}8-\mu_1}{5}\right) = 0{,}86433$, ce qui donne
$\dfrac{18{,}8-\mu_1}{5} = 1{,}10$.

\[
\mu_1 = 18{,}8 - 5\times 1{,}10 = 18{,}8 - 5{,}5 = 13{,}3.
\]
\end{solution}
\vspace{18cm}
$\mu_1=$ \boite{3}{1.5}

\pageblanche
\question[10] On demande à 3600 Suisses s'ils préfèrent la fondue ou la raclette. On obtient que 1849 personnes préfèrent la fondue. Y a-t-il suffisamment de preuves pour dire que la majorité des Suisses préfèrent la fondue, au seuil approximatif $\alpha = 0,05$? \textbf{\textit{Écrivez toutes les étapes du test et justifiez vos réponses!}}

\textit{Utilisez l'approximation suivante: $\displaystyle\sqrt{\frac{\frac{1849}{3600}\left(1-\frac{1849}{3600}\right)}{3600}}\approx0,008$}
\begin{solution}
\textbf{Étape 1 - Hypothèses.} Soit $p$ la proportion de Suisses préférant la fondue.
\[
H_0: p=0{,}5 \quad \text{contre} \quad H_1: p > 0{,}5.
\]
\textbf{Étape 2 - Statistique de test.}
\[
z_0 = \frac{\hat{p}-p_0}{\sqrt{\hat{p}(1-\hat{p})/n}}
= \frac{\frac{1849}{3600}-0{,}5}{0{,}008}
= \frac{\frac{49}{3600}}{0{,}008} \approx 1{,}70.
\]
\textbf{Étape 3 - Région de rejet.} Au seuil approximatif $\alpha=0{,}05$ (test unilatéral droit), on rejette $H_0$ si $z_0 > z_{0{,}05} = 1{,}645$.

\textbf{Décision.} Comme $1{,}70 > 1{,}645$, on \textbf{rejette $H_0$}. Il y a suffisamment de preuves, au seuil 5\%, pour affirmer que la majorité des Suisses préfèrent la fondue.
\end{solution}
\pageblanche
\question[10] Vous travaillez dans une usine qui produit des boulons, et le responsable de la qualité prétend que la variance de la longueur des boulons est de $4$ mm$^2$. Vous avez des raisons de penser que la variance est en réalité plus élevée et décidez de mener un test d'hypothèse.



Vous prenez un échantillon aléatoire de $20$ boulons et mesurez leurs longueurs. La variance des longueurs de l'échantillon est de $5,5$ mm$^2$. 
\\
\\
Effectuez un test d'hypothèse de seuil $\alpha=1\%$ pour déterminer s'il existe suffisamment de preuves pour rejeter l'affirmation du responsable de la qualité. Vous pouvez supposer que les longueurs des boulons suivent une distribution normale. \textbf{\textit{Écrivez toutes les étapes du test et justifiez vos réponses!}}

\begin{solution}
\textbf{Étape 1 - Hypothèses.} Soit $\sigma^2$ la variance des longueurs des boulons.
\[
H_0: \sigma^2 = 4 \quad \text{contre} \quad H_1: \sigma^2 > 4.
\]
\textbf{Postulat:} les longueurs suivent une loi normale.

\textbf{Étape 2 - Statistique de test.}
\[
\chi^2_0 = \frac{(n-1)s^2}{\sigma_0^2} = \frac{19 \times 5{,}5}{4} = \frac{104{,}5}{4} = 26{,}125.
\]
\textbf{Étape 3 - Région de rejet.} Au seuil $\alpha=1\%$, on rejette $H_0$ si $\chi^2_0 > \chi^2_{0{,}01\,;\,19}$. D'après la table, $\chi^2_{0{,}01\,;\,19} \approx 36{,}191$.

\textbf{Décision.} Comme $26{,}125 < 36{,}191$, on \textbf{ne rejette pas $H_0$}. Il n'y a pas suffisamment de preuves pour réfuter l'affirmation du responsable de la qualité.
\end{solution}



\end{questions}
\end{document} 
